\section{Contexto e Motivação}
\label{sec:contexto}

O setor elétrico brasileiro é estruturado em torno de uma matriz de geração predominantemente hidráulica, cuja dependência dos ciclos hidrológicos introduz incerteza intrínseca na oferta de energia ao longo do ano \cite{epe2025ben}. Essa incerteza se manifesta diretamente no Preço de Liquidação das Diferenças (PLD), o preço do mercado de curto prazo calculado semanalmente pela Câmara de Comercialização de Energia Elétrica (CCEE) com base no Custo Marginal de Operação (CMO) do Sistema Interligado Nacional (SIN). O PLD apresenta comportamento estatístico assimétrico e de cauda pesada: períodos de seca elevada podem empurrar o preço ao teto regulatório, ao passo que anos hidrológicos favoráveis o comprimem ao piso. Para os agentes que operam no Ambiente de Contratação Livre (ACL), essa variabilidade é a principal fonte de risco financeiro operacional \cite{guerreiro2017}.

Nesse contexto, as \emph{comercializadoras de energia} --- empresas que compram e vendem energia por meio de contratos bilaterais, sem depender exclusivamente da operação de ativos próprios para cobrir suas posições --- enfrentam um problema de gestão de portfólio multiperiódico. A cada mês, a empresa decide quais contratos firmar para os meses seguintes, sem conhecer os PLDs futuros nem a produção física que seus ativos entregarão. As decisões tomadas em um mês modificam a posição contratual de todos os meses subsequentes enquanto os contratos estiverem vigentes, criando uma cadeia de interdependências temporais que não pode ser ignorada em um modelo de otimização.

A abordagem clássica de resolver esse tipo de problema pelo valor esperado --- substituindo as variáveis aleatórias por suas médias e otimizando um problema determinístico único --- é conhecida por subestimar o risco de posições deficitárias e por produzir políticas que se revelam subótimas quando os cenários adversos se materializam \cite{birge2011}. A Programação Estocástica Multiestágio (PEM) oferece o arcabouço matemático adequado para tratar essa estrutura de decisão sequencial sob incerteza: cada estágio representa um mês de planejamento, as variáveis aleatórias do PLD e da produção física são representadas por uma árvore de cenários, e as decisões de contratação devem respeitar a \emph{não antecipatividade} --- a empresa não pode utilizar informação sobre o futuro ao firmar contratos no presente.

A motivação central deste trabalho é investigar como resolver o problema de gestão de portfólio da comercializadora de forma computacionalmente tratável para horizontes e discretizações de cenários realistas. Duas abordagens de solução são confrontadas: a \emph{Formulação Determinística Equivalente} (DEQ), que representa toda a árvore de cenários em um único programa linear de grande porte, e a \emph{Programação Dinâmica Dual Estocástica} (PDDE), que decompõe o problema estágio a estágio e aproxima as funções de valor futuro por cortes de Benders sem construir a árvore explicitamente \cite{pereira1991}. A comparação entre as duas abordagens --- em termos de escalabilidade computacional e qualidade da política gerada --- é o fio condutor do trabalho.
